\chapter{Objetivos}

\section{Objetivo General}
Optimizar la gestión de nóminas de Café y Vivero Verde Pradera a través 
del desarrollo de un sistema web automatizado, con el objetivo de aumentar 
la precisión en el cálculo de salarios y mejorar la eficiencia operativa. 
Este sistema también garantizará la protección de los datos personales de 
los empleados. La implementación se realizará tras la conclusión de los 
cursos de Ingeniería en Sistemas I, II y III en la Universidad Nacional de 
Costa Rica, programada para el final del primer semestre de 2026.

\section{Objetivos Específicos}

\subsection{Objetivo Especifico 1}
\textbf{Desarrollar una interfaz de usuario amigable e intuitiva:} Crear una interfaz 
web que facilite a los usuarios autorizados la entrada y gestión de datos 
de nómina, mejorando la experiencia del usuario y asegurando el manejo 
eficiente y libre de errores de las operaciones de nómina.

\subsection{Objetivo Especifico 2}
\textbf{Implementar un módulo de seguridad robusto:} Establecer protocolos de 
seguridad para proteger los datos personales y financieros de los empleados, 
restringiendo el acceso a la información exclusivamente a usuarios autorizados 
y asegurando la integridad y confidencialidad de los datos.

\newpage

\subsection{Objetivo Especifico 3}
\textbf{Automatizar el proceso de cálculo de nómina:} Desarrollar 
funcionalidades automáticas para calcular los salarios y 
generar reportes obligatorios para la Caja Costarricense del 
Seguro Social, el Ministerio de Trabajo y Seguridad, así como los 
comprobantes de pago para los empleados. Este desarrollo busca reducir 
errores humanos y aumentar la eficiencia del proceso de nómina.

\subsection{Objetivo Especifico 4}
\textbf{Implementar un módulo de administración de eventos laborales:} Desarrollar 
un sistema que permita a los usuarios autorizados registrar y gestionar 
manualmente fechas de vacaciones, periodos de incapacidad, y periodos de 
lactancia, entre otros eventos laborales significativos. Este módulo ayudará 
a mantener la consistencia de los datos y facilitará el acceso y la organización 
de la información, apoyando así una gestión precisa y eficiente de los eventos 
laborales dentro de la empresa.

\subsection{Objetivo Especifico 5}
\textbf{Asegurar la escalabilidad del sistema:} Diseñar el sistema de manera que 
pueda adaptarse fácilmente a futuras expansiones, cambios en la legislación 
laboral o ajustes en las políticas internas de la empresa sin comprometer su 
operatividad o seguridad.

\subsection{Objetivo Especifico 6}
\textbf{Capacitar a los usuarios en el uso del sistema:} Proveer formación adecuada 
para los usuarios autorizados sobre cómo operar el nuevo sistema de gestión 
de nóminas, asegurando así su correcta y eficaz utilización desde el momento 
de su implementación.

\subsection{Objetivo Especifico 7}
\textbf{Realizar pruebas exhaustivas antes de la implementación final:} Validar la 
funcionalidad, usabilidad y seguridad del sistema mediante pruebas rigurosas 
con casos de uso reales para garantizar que todos los componentes funcionen 
según lo previsto y sin fallos.
